% Options for packages loaded elsewhere
\PassOptionsToPackage{unicode}{hyperref}
\PassOptionsToPackage{hyphens}{url}
\documentclass[
]{article}
\usepackage{xcolor}
\usepackage[margin=1in]{geometry}
\usepackage{amsmath,amssymb}
\setcounter{secnumdepth}{-\maxdimen} % remove section numbering
\usepackage{iftex}
\ifPDFTeX
  \usepackage[T1]{fontenc}
  \usepackage[utf8]{inputenc}
  \usepackage{textcomp} % provide euro and other symbols
\else % if luatex or xetex
  \usepackage{unicode-math} % this also loads fontspec
  \defaultfontfeatures{Scale=MatchLowercase}
  \defaultfontfeatures[\rmfamily]{Ligatures=TeX,Scale=1}
\fi
\usepackage{lmodern}
\ifPDFTeX\else
  % xetex/luatex font selection
\fi
% Use upquote if available, for straight quotes in verbatim environments
\IfFileExists{upquote.sty}{\usepackage{upquote}}{}
\IfFileExists{microtype.sty}{% use microtype if available
  \usepackage[]{microtype}
  \UseMicrotypeSet[protrusion]{basicmath} % disable protrusion for tt fonts
}{}
\makeatletter
\@ifundefined{KOMAClassName}{% if non-KOMA class
  \IfFileExists{parskip.sty}{%
    \usepackage{parskip}
  }{% else
    \setlength{\parindent}{0pt}
    \setlength{\parskip}{6pt plus 2pt minus 1pt}}
}{% if KOMA class
  \KOMAoptions{parskip=half}}
\makeatother
\usepackage{color}
\usepackage{fancyvrb}
\newcommand{\VerbBar}{|}
\newcommand{\VERB}{\Verb[commandchars=\\\{\}]}
\DefineVerbatimEnvironment{Highlighting}{Verbatim}{commandchars=\\\{\}}
% Add ',fontsize=\small' for more characters per line
\usepackage{framed}
\definecolor{shadecolor}{RGB}{248,248,248}
\newenvironment{Shaded}{\begin{snugshade}}{\end{snugshade}}
\newcommand{\AlertTok}[1]{\textcolor[rgb]{0.94,0.16,0.16}{#1}}
\newcommand{\AnnotationTok}[1]{\textcolor[rgb]{0.56,0.35,0.01}{\textbf{\textit{#1}}}}
\newcommand{\AttributeTok}[1]{\textcolor[rgb]{0.13,0.29,0.53}{#1}}
\newcommand{\BaseNTok}[1]{\textcolor[rgb]{0.00,0.00,0.81}{#1}}
\newcommand{\BuiltInTok}[1]{#1}
\newcommand{\CharTok}[1]{\textcolor[rgb]{0.31,0.60,0.02}{#1}}
\newcommand{\CommentTok}[1]{\textcolor[rgb]{0.56,0.35,0.01}{\textit{#1}}}
\newcommand{\CommentVarTok}[1]{\textcolor[rgb]{0.56,0.35,0.01}{\textbf{\textit{#1}}}}
\newcommand{\ConstantTok}[1]{\textcolor[rgb]{0.56,0.35,0.01}{#1}}
\newcommand{\ControlFlowTok}[1]{\textcolor[rgb]{0.13,0.29,0.53}{\textbf{#1}}}
\newcommand{\DataTypeTok}[1]{\textcolor[rgb]{0.13,0.29,0.53}{#1}}
\newcommand{\DecValTok}[1]{\textcolor[rgb]{0.00,0.00,0.81}{#1}}
\newcommand{\DocumentationTok}[1]{\textcolor[rgb]{0.56,0.35,0.01}{\textbf{\textit{#1}}}}
\newcommand{\ErrorTok}[1]{\textcolor[rgb]{0.64,0.00,0.00}{\textbf{#1}}}
\newcommand{\ExtensionTok}[1]{#1}
\newcommand{\FloatTok}[1]{\textcolor[rgb]{0.00,0.00,0.81}{#1}}
\newcommand{\FunctionTok}[1]{\textcolor[rgb]{0.13,0.29,0.53}{\textbf{#1}}}
\newcommand{\ImportTok}[1]{#1}
\newcommand{\InformationTok}[1]{\textcolor[rgb]{0.56,0.35,0.01}{\textbf{\textit{#1}}}}
\newcommand{\KeywordTok}[1]{\textcolor[rgb]{0.13,0.29,0.53}{\textbf{#1}}}
\newcommand{\NormalTok}[1]{#1}
\newcommand{\OperatorTok}[1]{\textcolor[rgb]{0.81,0.36,0.00}{\textbf{#1}}}
\newcommand{\OtherTok}[1]{\textcolor[rgb]{0.56,0.35,0.01}{#1}}
\newcommand{\PreprocessorTok}[1]{\textcolor[rgb]{0.56,0.35,0.01}{\textit{#1}}}
\newcommand{\RegionMarkerTok}[1]{#1}
\newcommand{\SpecialCharTok}[1]{\textcolor[rgb]{0.81,0.36,0.00}{\textbf{#1}}}
\newcommand{\SpecialStringTok}[1]{\textcolor[rgb]{0.31,0.60,0.02}{#1}}
\newcommand{\StringTok}[1]{\textcolor[rgb]{0.31,0.60,0.02}{#1}}
\newcommand{\VariableTok}[1]{\textcolor[rgb]{0.00,0.00,0.00}{#1}}
\newcommand{\VerbatimStringTok}[1]{\textcolor[rgb]{0.31,0.60,0.02}{#1}}
\newcommand{\WarningTok}[1]{\textcolor[rgb]{0.56,0.35,0.01}{\textbf{\textit{#1}}}}
\usepackage{graphicx}
\makeatletter
\newsavebox\pandoc@box
\newcommand*\pandocbounded[1]{% scales image to fit in text height/width
  \sbox\pandoc@box{#1}%
  \Gscale@div\@tempa{\textheight}{\dimexpr\ht\pandoc@box+\dp\pandoc@box\relax}%
  \Gscale@div\@tempb{\linewidth}{\wd\pandoc@box}%
  \ifdim\@tempb\p@<\@tempa\p@\let\@tempa\@tempb\fi% select the smaller of both
  \ifdim\@tempa\p@<\p@\scalebox{\@tempa}{\usebox\pandoc@box}%
  \else\usebox{\pandoc@box}%
  \fi%
}
% Set default figure placement to htbp
\def\fps@figure{htbp}
\makeatother
\setlength{\emergencystretch}{3em} % prevent overfull lines
\providecommand{\tightlist}{%
  \setlength{\itemsep}{0pt}\setlength{\parskip}{0pt}}
\usepackage{bookmark}
\IfFileExists{xurl.sty}{\usepackage{xurl}}{} % add URL line breaks if available
\urlstyle{same}
\hypersetup{
  pdftitle={Swiss Synthetic-Textile Microfibre --- Monthly Time-Series Study},
  pdfauthor={Sustainability Analytics},
  hidelinks,
  pdfcreator={LaTeX via pandoc}}

\title{Swiss Synthetic-Textile Microfibre --- Monthly Time-Series Study}
\usepackage{etoolbox}
\makeatletter
\providecommand{\subtitle}[1]{% add subtitle to \maketitle
  \apptocmd{\@title}{\par {\large #1 \par}}{}{}
}
\makeatother
\subtitle{168 months of customs data · seasonality · SARIMA/SARIMAX ·
rolling-origin backtest}
\author{Sustainability Analytics}
\date{26 August 2026}

\begin{document}
\maketitle

{
\setcounter{tocdepth}{2}
\tableofcontents
}
\section{Why this study is different}\label{why-this-study-is-different}

Every earlier version of this analysis used \textbf{14 annual numbers}
--- a pre-aggregated summary. We have now gone back to the \textbf{raw
Swiss customs file (791,478 records)} and rebuilt the series at its
native \textbf{monthly} resolution: \textbf{168 months, January 2012 --
December 2025.} That one change lifts the analysis from ``a short line
you can barely forecast'' to a proper seasonal time series with real
statistical power.

\textbf{What monthly data unlocks.} \textbf{(1)} A genuine
\textbf{seasonal pattern} (strength ≈ 0.65) --- invisible in annual
data. \textbf{(2)} A real \textbf{SARIMA} model instead of a random
walk. \textbf{(3)} ACF/PACF with enough points to actually read.
\textbf{(4)} Shocks dated to the \textbf{month}, not the year.
\textbf{(5)} A seasonal forecast that still rolls up to the same annual
footprint.

\section{1. The monthly series}\label{the-monthly-series}

\begin{Shaded}
\begin{Highlighting}[]
\NormalTok{dm}\SpecialCharTok{$}\NormalTok{covid }\OtherTok{\textless{}{-}} \FunctionTok{as.integer}\NormalTok{(dm}\SpecialCharTok{$}\NormalTok{date }\SpecialCharTok{\textgreater{}=} \FunctionTok{as.Date}\NormalTok{(}\StringTok{"2020{-}03{-}01"}\NormalTok{) }\SpecialCharTok{\&}\NormalTok{ dm}\SpecialCharTok{$}\NormalTok{date }\SpecialCharTok{\textless{}=} \FunctionTok{as.Date}\NormalTok{(}\StringTok{"2021{-}06{-}01"}\NormalTok{))}
\FunctionTok{ggplot}\NormalTok{(dm, }\FunctionTok{aes}\NormalTok{(date, microfibre\_t))}\SpecialCharTok{+}
  \FunctionTok{annotate}\NormalTok{(}\StringTok{"rect"}\NormalTok{, }\AttributeTok{xmin=}\FunctionTok{as.Date}\NormalTok{(}\StringTok{"2020{-}03{-}01"}\NormalTok{), }\AttributeTok{xmax=}\FunctionTok{as.Date}\NormalTok{(}\StringTok{"2021{-}06{-}01"}\NormalTok{), }\AttributeTok{ymin=}\SpecialCharTok{{-}}\ConstantTok{Inf}\NormalTok{, }\AttributeTok{ymax=}\ConstantTok{Inf}\NormalTok{, }\AttributeTok{fill=}\NormalTok{POLY, }\AttributeTok{alpha=}\FloatTok{0.08}\NormalTok{)}\SpecialCharTok{+}
  \FunctionTok{geom\_line}\NormalTok{(}\AttributeTok{colour=}\NormalTok{TEAL, }\AttributeTok{linewidth=}\FloatTok{0.7}\NormalTok{)}\SpecialCharTok{+}
  \FunctionTok{geom\_smooth}\NormalTok{(}\AttributeTok{method=}\StringTok{"loess"}\NormalTok{, }\AttributeTok{span=}\FloatTok{0.3}\NormalTok{, }\AttributeTok{se=}\ConstantTok{FALSE}\NormalTok{, }\AttributeTok{colour=}\NormalTok{INK, }\AttributeTok{linewidth=}\FloatTok{0.9}\NormalTok{)}\SpecialCharTok{+}
  \FunctionTok{annotate}\NormalTok{(}\StringTok{"text"}\NormalTok{, }\AttributeTok{x=}\FunctionTok{as.Date}\NormalTok{(}\StringTok{"2020{-}10{-}01"}\NormalTok{), }\AttributeTok{y=}\FunctionTok{max}\NormalTok{(dm}\SpecialCharTok{$}\NormalTok{microfibre\_t), }\AttributeTok{label=}\StringTok{"COVID"}\NormalTok{, }\AttributeTok{colour=}\NormalTok{POLY, }\AttributeTok{size=}\FloatTok{3.6}\NormalTok{)}\SpecialCharTok{+}
  \FunctionTok{labs}\NormalTok{(}\AttributeTok{title=}\StringTok{"Monthly synthetic{-}microfibre release, 2012 to 2025"}\NormalTok{,}
       \AttributeTok{subtitle=}\StringTok{"168 monthly points (teal) with a smooth trend (black); the seasonal wave is visible to the eye"}\NormalTok{,}
       \AttributeTok{x=}\ConstantTok{NULL}\NormalTok{, }\AttributeTok{y=}\StringTok{"Microfibre released (t/month)"}\NormalTok{)}\SpecialCharTok{+}\FunctionTok{th}\NormalTok{()}
\end{Highlighting}
\end{Shaded}

\pandocbounded{\includegraphics[keepaspectratio]{time_series_monthly_files/figure-latex/series-1.pdf}}

\begin{Shaded}
\begin{Highlighting}[]
\FunctionTok{cat}\NormalTok{(}\FunctionTok{sprintf}\NormalTok{(}\StringTok{"Monthly mean \%.2f t | min \%.2f | max \%.2f | annual roll{-}up \textasciitilde{}\%.0f t/yr}\SpecialCharTok{\textbackslash{}n}\StringTok{"}\NormalTok{,}
            \FunctionTok{mean}\NormalTok{(y), }\FunctionTok{min}\NormalTok{(y), }\FunctionTok{max}\NormalTok{(y), }\FunctionTok{mean}\NormalTok{(y)}\SpecialCharTok{*}\DecValTok{12}\NormalTok{))}
\end{Highlighting}
\end{Shaded}

\begin{verbatim}
## Monthly mean 7.79 t | min 5.49 | max 9.88 | annual roll-up ~94 t/yr
\end{verbatim}

\section{2. Decomposition --- trend, season,
remainder}\label{decomposition-trend-season-remainder}

With monthly data an STL decomposition finally has a \textbf{seasonal}
component to find.

\begin{Shaded}
\begin{Highlighting}[]
\NormalTok{fit }\OtherTok{\textless{}{-}} \FunctionTok{stl}\NormalTok{(y, }\AttributeTok{s.window=}\StringTok{"periodic"}\NormalTok{, }\AttributeTok{robust=}\ConstantTok{TRUE}\NormalTok{)}
\FunctionTok{autoplot}\NormalTok{(fit) }\SpecialCharTok{+} \FunctionTok{labs}\NormalTok{(}\AttributeTok{title=}\StringTok{"STL decomposition: trend + seasonal + remainder"}\NormalTok{) }\SpecialCharTok{+} \FunctionTok{th}\NormalTok{()}
\end{Highlighting}
\end{Shaded}

\pandocbounded{\includegraphics[keepaspectratio]{time_series_monthly_files/figure-latex/stl-1.pdf}}

\begin{Shaded}
\begin{Highlighting}[]
\NormalTok{cm }\OtherTok{\textless{}{-}}\NormalTok{ fit}\SpecialCharTok{$}\NormalTok{time.series}
\NormalTok{Fs }\OtherTok{\textless{}{-}} \FunctionTok{round}\NormalTok{(}\FunctionTok{max}\NormalTok{(}\DecValTok{0}\NormalTok{,}\DecValTok{1}\SpecialCharTok{{-}}\FunctionTok{var}\NormalTok{(cm[,}\StringTok{"remainder"}\NormalTok{])}\SpecialCharTok{/}\FunctionTok{var}\NormalTok{(cm[,}\StringTok{"seasonal"}\NormalTok{]}\SpecialCharTok{+}\NormalTok{cm[,}\StringTok{"remainder"}\NormalTok{])),}\DecValTok{2}\NormalTok{)}
\NormalTok{Ft }\OtherTok{\textless{}{-}} \FunctionTok{round}\NormalTok{(}\FunctionTok{max}\NormalTok{(}\DecValTok{0}\NormalTok{,}\DecValTok{1}\SpecialCharTok{{-}}\FunctionTok{var}\NormalTok{(cm[,}\StringTok{"remainder"}\NormalTok{])}\SpecialCharTok{/}\FunctionTok{var}\NormalTok{(cm[,}\StringTok{"trend"}\NormalTok{]}\SpecialCharTok{+}\NormalTok{cm[,}\StringTok{"remainder"}\NormalTok{])),}\DecValTok{2}\NormalTok{)}
\FunctionTok{cat}\NormalTok{(}\FunctionTok{sprintf}\NormalTok{(}\StringTok{"Seasonal strength = \%.2f  |  Trend strength = \%.2f}\SpecialCharTok{\textbackslash{}n}\StringTok{"}\NormalTok{, Fs, Ft))}
\end{Highlighting}
\end{Shaded}

\begin{verbatim}
## Seasonal strength = 0.64  |  Trend strength = 0.45
\end{verbatim}

\textbf{Interpretation.} Seasonal strength ≈ \textbf{0.64} (0 = none, 1
= pure season) --- a \emph{real, moderate-to-strong} seasonal signal
that annual data completely hid. The trend is gentler (≈ 0.45),
confirming the earlier finding that the long-run rise is slow and
demand-led. The remainder is where genuine shocks live.

\section{3. The seasonal fingerprint}\label{the-seasonal-fingerprint}

\begin{Shaded}
\begin{Highlighting}[]
\NormalTok{seas }\OtherTok{\textless{}{-}} \FunctionTok{data.frame}\NormalTok{(}\AttributeTok{month=}\FunctionTok{factor}\NormalTok{(month.abb[}\FunctionTok{cycle}\NormalTok{(y)], }\AttributeTok{levels=}\NormalTok{month.abb), }\AttributeTok{s=}\FunctionTok{as.numeric}\NormalTok{(cm[,}\StringTok{"seasonal"}\NormalTok{]))}
\NormalTok{sidx }\OtherTok{\textless{}{-}}\NormalTok{ seas }\SpecialCharTok{\%\textgreater{}\%} \FunctionTok{group\_by}\NormalTok{(month) }\SpecialCharTok{\%\textgreater{}\%} \FunctionTok{summarise}\NormalTok{(}\AttributeTok{mean\_s=}\FunctionTok{mean}\NormalTok{(s), }\AttributeTok{.groups=}\StringTok{"drop"}\NormalTok{)}
\FunctionTok{ggplot}\NormalTok{(sidx, }\FunctionTok{aes}\NormalTok{(month, mean\_s, }\AttributeTok{fill=}\NormalTok{mean\_s}\SpecialCharTok{\textgreater{}}\DecValTok{0}\NormalTok{))}\SpecialCharTok{+}
  \FunctionTok{geom\_col}\NormalTok{()}\SpecialCharTok{+}\FunctionTok{scale\_fill\_manual}\NormalTok{(}\AttributeTok{values=}\FunctionTok{c}\NormalTok{(}\StringTok{\textasciigrave{}}\AttributeTok{TRUE}\StringTok{\textasciigrave{}}\OtherTok{=}\NormalTok{TEAL,}\StringTok{\textasciigrave{}}\AttributeTok{FALSE}\StringTok{\textasciigrave{}}\OtherTok{=}\NormalTok{POLY), }\AttributeTok{guide=}\StringTok{"none"}\NormalTok{)}\SpecialCharTok{+}
  \FunctionTok{labs}\NormalTok{(}\AttributeTok{title=}\StringTok{"Average seasonal effect by calendar month"}\NormalTok{,}
       \AttributeTok{subtitle=}\StringTok{"How much each month sits above/below the trend (tonnes); autumn peaks, December dips"}\NormalTok{,}
       \AttributeTok{x=}\ConstantTok{NULL}\NormalTok{, }\AttributeTok{y=}\StringTok{"Seasonal effect (t)"}\NormalTok{)}\SpecialCharTok{+}\FunctionTok{th}\NormalTok{()}
\end{Highlighting}
\end{Shaded}

\pandocbounded{\includegraphics[keepaspectratio]{time_series_monthly_files/figure-latex/season-1.pdf}}

\begin{Shaded}
\begin{Highlighting}[]
\NormalTok{top }\OtherTok{\textless{}{-}}\NormalTok{ sidx}\SpecialCharTok{$}\NormalTok{month[}\FunctionTok{which.max}\NormalTok{(sidx}\SpecialCharTok{$}\NormalTok{mean\_s)]; bot }\OtherTok{\textless{}{-}}\NormalTok{ sidx}\SpecialCharTok{$}\NormalTok{month[}\FunctionTok{which.min}\NormalTok{(sidx}\SpecialCharTok{$}\NormalTok{mean\_s)]}
\FunctionTok{cat}\NormalTok{(}\FunctionTok{sprintf}\NormalTok{(}\StringTok{"Peak month: \%s (+\%.2f t) | Trough month: \%s (\%.2f t)}\SpecialCharTok{\textbackslash{}n}\StringTok{"}\NormalTok{,}
\NormalTok{            top, }\FunctionTok{max}\NormalTok{(sidx}\SpecialCharTok{$}\NormalTok{mean\_s), bot, }\FunctionTok{min}\NormalTok{(sidx}\SpecialCharTok{$}\NormalTok{mean\_s)))}
\end{Highlighting}
\end{Shaded}

\begin{verbatim}
## Peak month: Oct (+0.79 t) | Trough month: Dec (-1.15 t)
\end{verbatim}

\textbf{Interpretation for the client.} Imports --- and therefore
microfibre --- \textbf{peak in autumn} (restocking for the autumn/winter
retail season) and \textbf{trough in December} (shelves already stocked;
customs slows over the holidays). For a brand this matters
operationally: the pollution load is \emph{not} flat across the year, so
interventions, reporting, and any wash-related messaging can be timed to
the peak months.

\section{4. Stationarity \& the seasonal
ACF/PACF}\label{stationarity-the-seasonal-acfpacf}

\begin{Shaded}
\begin{Highlighting}[]
\NormalTok{adf\_p}\OtherTok{\textless{}{-}}\ControlFlowTok{function}\NormalTok{(x) }\FunctionTok{suppressWarnings}\NormalTok{(}\FunctionTok{adf.test}\NormalTok{(x))}\SpecialCharTok{$}\NormalTok{p.value}
\NormalTok{kpss\_p}\OtherTok{\textless{}{-}}\ControlFlowTok{function}\NormalTok{(x) }\FunctionTok{suppressWarnings}\NormalTok{(}\FunctionTok{kpss.test}\NormalTok{(x))}\SpecialCharTok{$}\NormalTok{p.value}
\FunctionTok{data.frame}\NormalTok{(}
  \AttributeTok{Series=}\FunctionTok{c}\NormalTok{(}\StringTok{"Level"}\NormalTok{,}\StringTok{"Seasonal difference (lag 12)"}\NormalTok{,}\StringTok{"Seasonal + 1st difference"}\NormalTok{),}
  \AttributeTok{ADF\_p =}\FunctionTok{round}\NormalTok{(}\FunctionTok{c}\NormalTok{(}\FunctionTok{adf\_p}\NormalTok{(y), }\FunctionTok{adf\_p}\NormalTok{(}\FunctionTok{diff}\NormalTok{(y,}\DecValTok{12}\NormalTok{)), }\FunctionTok{adf\_p}\NormalTok{(}\FunctionTok{diff}\NormalTok{(}\FunctionTok{diff}\NormalTok{(y,}\DecValTok{12}\NormalTok{)))),}\DecValTok{3}\NormalTok{),}
  \AttributeTok{KPSS\_p=}\FunctionTok{round}\NormalTok{(}\FunctionTok{c}\NormalTok{(}\FunctionTok{kpss\_p}\NormalTok{(y), }\FunctionTok{kpss\_p}\NormalTok{(}\FunctionTok{diff}\NormalTok{(y,}\DecValTok{12}\NormalTok{)), }\FunctionTok{kpss\_p}\NormalTok{(}\FunctionTok{diff}\NormalTok{(}\FunctionTok{diff}\NormalTok{(y,}\DecValTok{12}\NormalTok{)))),}\DecValTok{3}\NormalTok{))}
\end{Highlighting}
\end{Shaded}

\begin{verbatim}
##                         Series ADF_p KPSS_p
## 1                        Level 0.048   0.01
## 2 Seasonal difference (lag 12) 0.231   0.10
## 3    Seasonal + 1st difference 0.010   0.10
\end{verbatim}

\begin{Shaded}
\begin{Highlighting}[]
\FunctionTok{par}\NormalTok{(}\AttributeTok{mfrow=}\FunctionTok{c}\NormalTok{(}\DecValTok{1}\NormalTok{,}\DecValTok{2}\NormalTok{), }\AttributeTok{col.axis=}\NormalTok{MUT, }\AttributeTok{col.lab=}\NormalTok{MUT, }\AttributeTok{col.main=}\NormalTok{INK, }\AttributeTok{mar=}\FunctionTok{c}\NormalTok{(}\DecValTok{4}\NormalTok{,}\DecValTok{4}\NormalTok{,}\DecValTok{3}\NormalTok{,}\DecValTok{1}\NormalTok{))}
\FunctionTok{Acf}\NormalTok{(y, }\AttributeTok{lag.max=}\DecValTok{36}\NormalTok{, }\AttributeTok{main=}\StringTok{"ACF — level (note spikes every 12 months)"}\NormalTok{)}
\FunctionTok{Pacf}\NormalTok{(y, }\AttributeTok{lag.max=}\DecValTok{36}\NormalTok{, }\AttributeTok{main=}\StringTok{"PACF — level"}\NormalTok{)}
\end{Highlighting}
\end{Shaded}

\pandocbounded{\includegraphics[keepaspectratio]{time_series_monthly_files/figure-latex/acf-1.pdf}}

\textbf{Interpretation.} Now that we have 168 points the ACF is
\textbf{readable} and shows clear \textbf{spikes at lags 12, 24} --- the
statistical signature of annual seasonality (utterly absent from the
14-point annual ACF). A \textbf{seasonal difference} (subtracting the
same month a year earlier) is what makes the series stationary --- which
is exactly the \texttt{(…)(0,1,1){[}12{]}} term the model selects next.

\section{5. Models \& the honest
backtest}\label{models-the-honest-backtest}

We compare three forecasters and score them by \textbf{MASE} (below 1
beats the seasonal-naïve benchmark):

\begin{itemize}
\tightlist
\item
  \textbf{Seasonal naïve} --- ``next January = last January.'' The
  benchmark.
\item
  \textbf{SARIMA} --- seasonal ARIMA chosen automatically.
\item
  \textbf{SARIMAX} --- the same, plus a \textbf{COVID event regressor}
  (an exogenous variable).
\end{itemize}

\begin{Shaded}
\begin{Highlighting}[]
\NormalTok{xreg }\OtherTok{\textless{}{-}} \FunctionTok{matrix}\NormalTok{(dm}\SpecialCharTok{$}\NormalTok{covid, }\AttributeTok{ncol=}\DecValTok{1}\NormalTok{, }\AttributeTok{dimnames=}\FunctionTok{list}\NormalTok{(}\ConstantTok{NULL}\NormalTok{,}\StringTok{"covid"}\NormalTok{))}
\NormalTok{m\_sarima  }\OtherTok{\textless{}{-}} \FunctionTok{auto.arima}\NormalTok{(y)                    }\CommentTok{\# search the seasonal order ONCE}
\NormalTok{m\_sarimax }\OtherTok{\textless{}{-}} \FunctionTok{auto.arima}\NormalTok{(y, }\AttributeTok{xreg=}\NormalTok{xreg)}
\NormalTok{ord  }\OtherTok{\textless{}{-}} \FunctionTok{arimaorder}\NormalTok{(m\_sarima);  ordX }\OtherTok{\textless{}{-}} \FunctionTok{arimaorder}\NormalTok{(m\_sarimax)}
\FunctionTok{cat}\NormalTok{(}\StringTok{"SARIMA :"}\NormalTok{, }\FunctionTok{sprintf}\NormalTok{(}\StringTok{"ARIMA(\%d,\%d,\%d)(\%d,\%d,\%d)[12]"}\NormalTok{, ord[}\DecValTok{1}\NormalTok{],ord[}\DecValTok{2}\NormalTok{],ord[}\DecValTok{3}\NormalTok{],ord[}\DecValTok{4}\NormalTok{],ord[}\DecValTok{5}\NormalTok{],ord[}\DecValTok{6}\NormalTok{]), }\StringTok{"}\SpecialCharTok{\textbackslash{}n}\StringTok{"}\NormalTok{)}
\end{Highlighting}
\end{Shaded}

\begin{verbatim}
## SARIMA : ARIMA(0,0,3)(0,1,1)[12]
\end{verbatim}

\begin{Shaded}
\begin{Highlighting}[]
\CommentTok{\# Refit the CHOSEN orders at each origin (fast) rather than re{-}searching every time.}
\NormalTok{h\_test }\OtherTok{\textless{}{-}} \DecValTok{24}\NormalTok{; origins }\OtherTok{\textless{}{-}}\NormalTok{ (n}\SpecialCharTok{{-}}\NormalTok{h\_test)}\SpecialCharTok{:}\NormalTok{(n}\DecValTok{{-}1}\NormalTok{)}
\NormalTok{err }\OtherTok{\textless{}{-}} \FunctionTok{list}\NormalTok{(}\AttributeTok{SeasonalNaive=}\FunctionTok{numeric}\NormalTok{(), }\AttributeTok{SARIMA=}\FunctionTok{numeric}\NormalTok{(), }\AttributeTok{SARIMAX=}\FunctionTok{numeric}\NormalTok{())}
\NormalTok{safefc }\OtherTok{\textless{}{-}} \ControlFlowTok{function}\NormalTok{(train, order, seasonal, }\AttributeTok{xr=}\ConstantTok{NULL}\NormalTok{, }\AttributeTok{xf=}\ConstantTok{NULL}\NormalTok{)\{}
  \FunctionTok{tryCatch}\NormalTok{(\{}
\NormalTok{    fit }\OtherTok{\textless{}{-}} \ControlFlowTok{if}\NormalTok{ (}\FunctionTok{is.null}\NormalTok{(xr)) }\FunctionTok{Arima}\NormalTok{(train, }\AttributeTok{order=}\NormalTok{order, }\AttributeTok{seasonal=}\NormalTok{seasonal)}
           \ControlFlowTok{else}             \FunctionTok{Arima}\NormalTok{(train, }\AttributeTok{order=}\NormalTok{order, }\AttributeTok{seasonal=}\NormalTok{seasonal, }\AttributeTok{xreg=}\NormalTok{xr)}
    \ControlFlowTok{if}\NormalTok{ (}\FunctionTok{is.null}\NormalTok{(xf)) }\FunctionTok{as.numeric}\NormalTok{(}\FunctionTok{forecast}\NormalTok{(fit, }\AttributeTok{h=}\DecValTok{1}\NormalTok{)}\SpecialCharTok{$}\NormalTok{mean)}
    \ControlFlowTok{else}             \FunctionTok{as.numeric}\NormalTok{(}\FunctionTok{forecast}\NormalTok{(fit, }\AttributeTok{h=}\DecValTok{1}\NormalTok{, }\AttributeTok{xreg=}\NormalTok{xf)}\SpecialCharTok{$}\NormalTok{mean)}
\NormalTok{  \}, }\AttributeTok{error=}\ControlFlowTok{function}\NormalTok{(e) }\FunctionTok{as.numeric}\NormalTok{(}\FunctionTok{forecast}\NormalTok{(}\FunctionTok{snaive}\NormalTok{(train), }\AttributeTok{h=}\DecValTok{1}\NormalTok{)}\SpecialCharTok{$}\NormalTok{mean))}
\NormalTok{\}}
\ControlFlowTok{for}\NormalTok{ (t }\ControlFlowTok{in}\NormalTok{ origins)\{}
\NormalTok{  ytr }\OtherTok{\textless{}{-}} \FunctionTok{ts}\NormalTok{(y[}\DecValTok{1}\SpecialCharTok{:}\NormalTok{t], }\AttributeTok{start=}\FunctionTok{c}\NormalTok{(}\DecValTok{2012}\NormalTok{,}\DecValTok{1}\NormalTok{), }\AttributeTok{frequency=}\DecValTok{12}\NormalTok{); actual }\OtherTok{\textless{}{-}}\NormalTok{ y[t}\SpecialCharTok{+}\DecValTok{1}\NormalTok{]}
\NormalTok{  err}\SpecialCharTok{$}\NormalTok{SeasonalNaive }\OtherTok{\textless{}{-}} \FunctionTok{c}\NormalTok{(err}\SpecialCharTok{$}\NormalTok{SeasonalNaive, actual }\SpecialCharTok{{-}} \FunctionTok{as.numeric}\NormalTok{(}\FunctionTok{forecast}\NormalTok{(}\FunctionTok{snaive}\NormalTok{(ytr), }\AttributeTok{h=}\DecValTok{1}\NormalTok{)}\SpecialCharTok{$}\NormalTok{mean))}
\NormalTok{  err}\SpecialCharTok{$}\NormalTok{SARIMA }\OtherTok{\textless{}{-}} \FunctionTok{c}\NormalTok{(err}\SpecialCharTok{$}\NormalTok{SARIMA, actual }\SpecialCharTok{{-}} \FunctionTok{safefc}\NormalTok{(ytr, ord[}\DecValTok{1}\SpecialCharTok{:}\DecValTok{3}\NormalTok{], ord[}\DecValTok{4}\SpecialCharTok{:}\DecValTok{6}\NormalTok{]))}
\NormalTok{  xr}\OtherTok{\textless{}{-}}\FunctionTok{matrix}\NormalTok{(xreg[}\DecValTok{1}\SpecialCharTok{:}\NormalTok{t,],}\AttributeTok{ncol=}\DecValTok{1}\NormalTok{,}\AttributeTok{dimnames=}\FunctionTok{list}\NormalTok{(}\ConstantTok{NULL}\NormalTok{,}\StringTok{"covid"}\NormalTok{)); xf}\OtherTok{\textless{}{-}}\FunctionTok{matrix}\NormalTok{(xreg[t}\SpecialCharTok{+}\DecValTok{1}\NormalTok{,],}\AttributeTok{ncol=}\DecValTok{1}\NormalTok{,}\AttributeTok{dimnames=}\FunctionTok{list}\NormalTok{(}\ConstantTok{NULL}\NormalTok{,}\StringTok{"covid"}\NormalTok{))}
\NormalTok{  err}\SpecialCharTok{$}\NormalTok{SARIMAX }\OtherTok{\textless{}{-}} \FunctionTok{c}\NormalTok{(err}\SpecialCharTok{$}\NormalTok{SARIMAX, actual }\SpecialCharTok{{-}} \FunctionTok{safefc}\NormalTok{(ytr, ordX[}\DecValTok{1}\SpecialCharTok{:}\DecValTok{3}\NormalTok{], ordX[}\DecValTok{4}\SpecialCharTok{:}\DecValTok{6}\NormalTok{], xr, xf))}
\NormalTok{\}}
\NormalTok{naive\_mae }\OtherTok{\textless{}{-}} \FunctionTok{mean}\NormalTok{(}\FunctionTok{abs}\NormalTok{(err}\SpecialCharTok{$}\NormalTok{SeasonalNaive))}
\NormalTok{bt }\OtherTok{\textless{}{-}} \FunctionTok{lapply}\NormalTok{(}\FunctionTok{names}\NormalTok{(err), }\ControlFlowTok{function}\NormalTok{(m) }\FunctionTok{data.frame}\NormalTok{(}\AttributeTok{Model=}\NormalTok{m,}
  \AttributeTok{RMSE=}\FunctionTok{round}\NormalTok{(}\FunctionTok{sqrt}\NormalTok{(}\FunctionTok{mean}\NormalTok{(err[[m]]}\SpecialCharTok{\^{}}\DecValTok{2}\NormalTok{)),}\DecValTok{3}\NormalTok{), }\AttributeTok{MAE=}\FunctionTok{round}\NormalTok{(}\FunctionTok{mean}\NormalTok{(}\FunctionTok{abs}\NormalTok{(err[[m]])),}\DecValTok{3}\NormalTok{),}
  \AttributeTok{MASE=}\FunctionTok{round}\NormalTok{(}\FunctionTok{mean}\NormalTok{(}\FunctionTok{abs}\NormalTok{(err[[m]]))}\SpecialCharTok{/}\NormalTok{naive\_mae,}\DecValTok{3}\NormalTok{))) }\SpecialCharTok{\%\textgreater{}\%} \FunctionTok{bind\_rows}\NormalTok{() }\SpecialCharTok{\%\textgreater{}\%} \FunctionTok{arrange}\NormalTok{(MASE)}
\NormalTok{bt}
\end{Highlighting}
\end{Shaded}

\begin{verbatim}
##           Model  RMSE   MAE  MASE
## 1       SARIMAX 0.517 0.378 0.617
## 2        SARIMA 0.532 0.399 0.652
## 3 SeasonalNaive 0.726 0.612 1.000
\end{verbatim}

\begin{Shaded}
\begin{Highlighting}[]
\NormalTok{best }\OtherTok{\textless{}{-}}\NormalTok{ bt}\SpecialCharTok{$}\NormalTok{Model[}\DecValTok{1}\NormalTok{]}
\FunctionTok{cat}\NormalTok{(}\StringTok{"Best by backtest MASE:"}\NormalTok{, best, }\StringTok{"}\SpecialCharTok{\textbackslash{}n}\StringTok{"}\NormalTok{)}
\end{Highlighting}
\end{Shaded}

\begin{verbatim}
## Best by backtest MASE: SARIMAX
\end{verbatim}

\textbf{Interpretation.} The seasonal models are \textbf{real} models
now --- ARIMA(0,0,3)(0,1,1){[}12{]} --- with a seasonal-difference and
seasonal-MA term the annual data could never support. On the rolling
backtest the best model is \textbf{SARIMAX} (MASE 0.617); a MASE below 1
means it genuinely beats ``assume next month equals the same month last
year.'' \emph{(With a real seasonal signal, even the naïve benchmark is
now strong --- so beating it is a meaningful result, not a formality.)}

\section{6. Seasonal forecast to 2030}\label{seasonal-forecast-to-2030}

\begin{Shaded}
\begin{Highlighting}[]
\NormalTok{fc\_h }\OtherTok{\textless{}{-}} \DecValTok{60}  \CommentTok{\# months to Dec 2030}
\NormalTok{xf }\OtherTok{\textless{}{-}} \FunctionTok{matrix}\NormalTok{(}\DecValTok{0}\NormalTok{, }\AttributeTok{nrow=}\NormalTok{fc\_h, }\AttributeTok{ncol=}\DecValTok{1}\NormalTok{, }\AttributeTok{dimnames=}\FunctionTok{list}\NormalTok{(}\ConstantTok{NULL}\NormalTok{,}\StringTok{"covid"}\NormalTok{))}
\NormalTok{fc }\OtherTok{\textless{}{-}} \FunctionTok{forecast}\NormalTok{(m\_sarimax, }\AttributeTok{xreg=}\NormalTok{xf, }\AttributeTok{h=}\NormalTok{fc\_h)}
\FunctionTok{autoplot}\NormalTok{(fc)}\SpecialCharTok{+}\FunctionTok{labs}\NormalTok{(}\AttributeTok{title=}\StringTok{"SARIMAX seasonal forecast to 2030"}\NormalTok{,}
   \AttributeTok{subtitle=}\StringTok{"Monthly forecast keeps the seasonal wave; shaded = 80/95\% intervals"}\NormalTok{,}
   \AttributeTok{x=}\ConstantTok{NULL}\NormalTok{,}\AttributeTok{y=}\StringTok{"Microfibre (t/month)"}\NormalTok{)}\SpecialCharTok{+}\FunctionTok{th}\NormalTok{()}
\end{Highlighting}
\end{Shaded}

\pandocbounded{\includegraphics[keepaspectratio]{time_series_monthly_files/figure-latex/forecast-1.pdf}}

\begin{Shaded}
\begin{Highlighting}[]
\NormalTok{ann }\OtherTok{\textless{}{-}} \FunctionTok{data.frame}\NormalTok{(}\AttributeTok{year=}\DecValTok{2026}\SpecialCharTok{:}\DecValTok{2030}\NormalTok{,}
  \AttributeTok{t =} \FunctionTok{sapply}\NormalTok{(}\DecValTok{2026}\SpecialCharTok{:}\DecValTok{2030}\NormalTok{, }\ControlFlowTok{function}\NormalTok{(Y)\{ idx}\OtherTok{\textless{}{-}}\FunctionTok{which}\NormalTok{(}\FunctionTok{floor}\NormalTok{(}\FunctionTok{time}\NormalTok{(fc}\SpecialCharTok{$}\NormalTok{mean)}\SpecialCharTok{+}\FloatTok{1e{-}6}\NormalTok{)}\SpecialCharTok{==}\NormalTok{Y); }\FunctionTok{sum}\NormalTok{(fc}\SpecialCharTok{$}\NormalTok{mean[idx]) \}))}
\FunctionTok{cat}\NormalTok{(}\StringTok{"Annual roll{-}up of the monthly forecast (t/yr):}\SpecialCharTok{\textbackslash{}n}\StringTok{"}\NormalTok{); }\FunctionTok{print}\NormalTok{(ann)}
\end{Highlighting}
\end{Shaded}

\begin{verbatim}
## Annual roll-up of the monthly forecast (t/yr):
\end{verbatim}

\begin{verbatim}
##   year        t
## 1 2026 97.93735
## 2 2027 96.47276
## 3 2028 95.62552
## 4 2029 95.24477
## 5 2030 95.07672
\end{verbatim}

\textbf{Interpretation for the client.} The forecast now carries the
\textbf{seasonal wave forward} --- you can see \emph{when} in each year
the load lands, not just the annual total. Rolled up to years it stays
close to the \textbf{\textasciitilde90--100 t/yr} the annual study found
(a good consistency check), but the monthly view is what lets a brand
plan around peaks and gives far tighter, better-calibrated uncertainty.

\section{7. Did the fibre mix actually shift? (product-code
check)}\label{did-the-fibre-mix-actually-shift-product-code-check}

The whole project assumes a \textbf{flat 69\% synthetic share}. With the
raw file we can finally interrogate that directly --- the product codes
name the fibre for part of the trade. The honest result is a finding in
itself.

\begin{Shaded}
\begin{Highlighting}[]
\NormalTok{fb }\OtherTok{\textless{}{-}} \FunctionTok{read\_csv}\NormalTok{(}\FunctionTok{file.path}\NormalTok{(params}\SpecialCharTok{$}\NormalTok{data\_dir,}\StringTok{"fibre\_share\_by\_year.csv"}\NormalTok{), }\AttributeTok{show\_col\_types=}\ConstantTok{FALSE}\NormalTok{) }\SpecialCharTok{\%\textgreater{}\%} \FunctionTok{filter}\NormalTok{(year}\SpecialCharTok{\textless{}=}\DecValTok{2025}\NormalTok{)}
\NormalTok{fb }\OtherTok{\textless{}{-}}\NormalTok{ fb }\SpecialCharTok{\%\textgreater{}\%} \FunctionTok{mutate}\NormalTok{(}
  \AttributeTok{syn\_side =}\NormalTok{ kg\_syn }\SpecialCharTok{+}\NormalTok{ kg\_nonwoven,                      }\CommentTok{\# synthetic + (mostly{-}synthetic) nonwovens}
  \AttributeTok{explicit =}\NormalTok{ kg\_syn }\SpecialCharTok{+}\NormalTok{ kg\_nonwoven }\SpecialCharTok{+}\NormalTok{ kg\_nat,             }\CommentTok{\# fibre is named}
  \AttributeTok{total    =}\NormalTok{ kg\_syn }\SpecialCharTok{+}\NormalTok{ kg\_nonwoven }\SpecialCharTok{+}\NormalTok{ kg\_nat }\SpecialCharTok{+}\NormalTok{ kg\_glass }\SpecialCharTok{+}\NormalTok{ kg\_unspec,}
  \AttributeTok{syn\_share\_explicit =} \DecValTok{100}\SpecialCharTok{*}\NormalTok{syn\_side}\SpecialCharTok{/}\NormalTok{explicit,}
  \AttributeTok{coverage =} \DecValTok{100}\SpecialCharTok{*}\NormalTok{(explicit}\SpecialCharTok{+}\NormalTok{kg\_glass)}\SpecialCharTok{/}\NormalTok{total)}
\NormalTok{cov }\OtherTok{\textless{}{-}} \FunctionTok{round}\NormalTok{(}\FunctionTok{mean}\NormalTok{(fb}\SpecialCharTok{$}\NormalTok{coverage),}\DecValTok{0}\NormalTok{)}
\FunctionTok{ggplot}\NormalTok{(fb, }\FunctionTok{aes}\NormalTok{(year, syn\_share\_explicit))}\SpecialCharTok{+}
  \FunctionTok{geom\_hline}\NormalTok{(}\AttributeTok{yintercept=}\DecValTok{69}\NormalTok{, }\AttributeTok{linetype=}\StringTok{"dashed"}\NormalTok{, }\AttributeTok{colour=}\NormalTok{POLY)}\SpecialCharTok{+}
  \FunctionTok{annotate}\NormalTok{(}\StringTok{"text"}\NormalTok{, }\AttributeTok{x=}\FloatTok{2013.4}\NormalTok{, }\AttributeTok{y=}\FloatTok{70.6}\NormalTok{, }\AttributeTok{label=}\StringTok{"project assumption (69\%)"}\NormalTok{, }\AttributeTok{colour=}\NormalTok{POLY, }\AttributeTok{size=}\FloatTok{3.4}\NormalTok{)}\SpecialCharTok{+}
  \FunctionTok{geom\_line}\NormalTok{(}\AttributeTok{colour=}\NormalTok{TEAL, }\AttributeTok{linewidth=}\FloatTok{1.1}\NormalTok{)}\SpecialCharTok{+}\FunctionTok{geom\_point}\NormalTok{(}\AttributeTok{colour=}\NormalTok{INK, }\AttributeTok{size=}\FloatTok{1.8}\NormalTok{)}\SpecialCharTok{+}
  \FunctionTok{scale\_y\_continuous}\NormalTok{(}\AttributeTok{limits=}\FunctionTok{c}\NormalTok{(}\DecValTok{40}\NormalTok{,}\DecValTok{75}\NormalTok{))}\SpecialCharTok{+}
  \FunctionTok{labs}\NormalTok{(}\AttributeTok{title=}\StringTok{"Switzerland\textquotesingle{}s synthetic share, where the product codes name a fibre"}\NormalTok{,}
       \AttributeTok{subtitle=}\FunctionTok{sprintf}\NormalTok{(}\StringTok{"Only \textasciitilde{}\%d\%\% of imports specify a fibre; within that, the synthetic share is flat, not rising"}\NormalTok{, cov),}
       \AttributeTok{x=}\ConstantTok{NULL}\NormalTok{, }\AttributeTok{y=}\StringTok{"Synthetic share of fibre{-}explicit imports (\%)"}\NormalTok{)}\SpecialCharTok{+}\FunctionTok{th}\NormalTok{()}
\end{Highlighting}
\end{Shaded}

\pandocbounded{\includegraphics[keepaspectratio]{time_series_monthly_files/figure-latex/fibre-1.pdf}}

\begin{Shaded}
\begin{Highlighting}[]
\FunctionTok{cat}\NormalTok{(}\FunctionTok{sprintf}\NormalTok{(}\StringTok{"Fibre{-}explicit coverage: \textasciitilde{}\%d\%\% of kg | synthetic share \%.0f\%\% (2012) {-}\textgreater{} \%.0f\%\% (2025), essentially flat}\SpecialCharTok{\textbackslash{}n}\StringTok{"}\NormalTok{,}
\NormalTok{            cov, fb}\SpecialCharTok{$}\NormalTok{syn\_share\_explicit[}\DecValTok{1}\NormalTok{], }\FunctionTok{tail}\NormalTok{(fb}\SpecialCharTok{$}\NormalTok{syn\_share\_explicit,}\DecValTok{1}\NormalTok{)))}
\end{Highlighting}
\end{Shaded}

\begin{verbatim}
## Fibre-explicit coverage: ~30% of kg | synthetic share 58% (2012) -> 55% (2025), essentially flat
\end{verbatim}

\textbf{Interpretation for the client.} Two honest takeaways.
\textbf{First, coverage is limited:} only about \textbf{30\%} of imports
name their fibre in the customs codes --- the other \textasciitilde70\%
are classified by \emph{garment type} (a ``T-shirt'' or ``trousers''
carries no fibre in the data). So the customs file \textbf{cannot pin
the exact absolute synthetic share} without an HS-level concordance.
\textbf{Second --- and this is the useful part --- where fibre \emph{is}
named, the synthetic share is \textasciitilde55\% and essentially flat}
across 2012--2025, not rising. That \textbf{confirms two things we'd
assumed on faith:} the fibre mix is \textbf{not a hidden driver} of the
pollution rise (backing up the Kaya result that population, not
fibre-shift, carries it), and holding the share \emph{constant} in the
model is therefore \textbf{justified} --- the trend comes from volume,
not from Switzerland switching to synthetics. \emph{(The absolute level
--- 69\% vs the \textasciitilde55\% seen in the explicit subset ---
stays a calibration choice, since the unclassified 70\% is a garment mix
we can't resolve here; this is why the level, not the trend, remains the
model's main assumption.)}

\section{8. What this changes}\label{what-this-changes}

\textbf{Before (annual, 14 points):} a near-flat random walk, no
seasonality, ACF/PACF unreadable, forecast ≈ ``carry the level
forward.'' \textbf{Now (monthly, 168 points):} a genuine
\textbf{seasonal SARIMA/SARIMAX}, a measured \textbf{autumn-peak /
December-trough} cycle, readable diagnostics, a backtest against a
strong seasonal benchmark, and a monthly forecast that still reconciles
to the same annual footprint.

We also \textbf{interrogated the 69\% synthetic assumption} against the
product codes (§7): the share is flat, so fibre-mix is not driving the
trend --- the assumption is safe for the trend, even if its absolute
level can't be pinned from customs codes alone.

\textbf{Still open (worth naming honestly).} Pinning the \emph{absolute}
synthetic share would need an HS-level fibre concordance for the
\textasciitilde70\% of imports classified only by garment type. And a
\textbf{monthly temperature series} could now be added as a second
exogenous driver of washing behaviour, since we finally have a monthly
target to attach it to.

\begin{center}\rule{0.5\linewidth}{0.5pt}\end{center}

\emph{Reproducible: \texttt{data/textile\_imports\_by\_month.csv} and
\texttt{data/fibre\_share\_by\_year.csv} are aggregated from the raw
791k-record customs file; knit this file to regenerate every figure.}

\end{document}
